% ============================================================================
% CHAPTER 8: RLHF AND ALIGNMENT
% ============================================================================
% Covers the full alignment pipeline: reward modeling, PPO, DPO, CAI,
% process reward models, and the practical challenges of aligning LLMs.
% ============================================================================

\chapter{RLHF and Alignment}
\label{chap:rlhf_alignment}

\begin{tldr}
Alignment transforms capable but potentially harmful base models into
helpful, harmless, and honest assistants.  This chapter covers the complete
alignment pipeline: reward model training from human preferences, PPO
optimization with KL constraints, DPO as a simpler alternative, and
Constitutional AI for scalable oversight.  Understanding RLHF is essential
for interviews at companies building or deploying LLMs---alignment is what
turned GPT-3 into ChatGPT, and every frontier lab now invests heavily in
these techniques.
\end{tldr}

% ============================================================================
\section{The Alignment Problem}
\label{sec:alignment_problem}
% ============================================================================

A pretrained language model is a powerful next-token predictor, but it is
not a \emph{useful assistant}.  It will happily continue toxic text, produce
plausible-sounding misinformation, refuse to follow instructions, or ramble
incoherently---because its training objective (predict the next token in
internet text) has no notion of helpfulness, safety, or honesty.

\paragraph{The three H's.}
Anthropic's framing (Askell et al., 2021) crystallized the alignment target
into three properties:

\begin{enumerate}[leftmargin=*]
\item \textbf{Helpful:} The model follows user instructions, provides
  relevant information, and accomplishes requested tasks effectively.
\item \textbf{Harmless:} The model refuses dangerous requests, avoids
  generating toxic or biased content, and does not assist in harmful
  activities.
\item \textbf{Honest:} The model avoids fabricating information, expresses
  uncertainty when appropriate, and does not systematically deceive the
  user.
\end{enumerate}

\paragraph{Why supervised fine-tuning alone is insufficient.}
Supervised fine-tuning (SFT) on human-written demonstrations is a necessary
first step, but it has fundamental limitations:

\begin{itemize}[leftmargin=*]
\item \textbf{Sycophancy:} SFT models learn to produce ``human-like''
  responses, which includes agreeing with the user even when they are wrong.
\item \textbf{Hallucination:} Imitating human writing style does not teach
  the model what it does and does not know.
\item \textbf{Harmful compliance:} SFT models can be instruction-following
  to a fault, complying with harmful requests if phrased as instructions.
\item \textbf{Mode averaging:} Human demonstrations exhibit diverse styles.
  SFT averages over these modes, producing bland, uncommitted outputs rather
  than consistently high-quality ones.
\item \textbf{Comparison is easier than generation:} It is easier for a
  human to say ``response A is better than response B'' than to write a
  perfect response from scratch.  RLHF leverages this by learning from
  \emph{preferences} rather than demonstrations.
\end{itemize}

\begin{keyinsight}[Why Preferences Beat Demonstrations]
The core insight of RLHF is that human \emph{preferences} contain richer
signal than human \emph{demonstrations}.  Writing a perfect response
requires solving the task; ranking two responses only requires recognizing
quality.  Humans are better judges than performers---this is true for
writing, coding, and virtually all creative tasks.  RLHF exploits this
asymmetry by training a reward model on easy-to-collect preference data,
then optimizing the LLM against that reward model.
\end{keyinsight}

\subsection{The Full Alignment Pipeline}

The standard alignment pipeline, introduced by InstructGPT (Ouyang et al.,
2022), consists of three stages.

\begin{figure}[H]
\centering
\begin{tikzpicture}[
    node distance=0.9cm and 0.7cm,
    stage/.style={rectangle, draw=deepblue, fill=lightblue, thick,
                  minimum height=1.1cm, minimum width=2.6cm,
                  font=\small\sffamily, rounded corners=3pt,
                  text width=2.4cm, align=center},
    datastage/.style={rectangle, draw=gray!60, fill=gray!8,
                      minimum height=0.7cm, minimum width=2.4cm,
                      font=\scriptsize\sffamily, rounded corners=2pt,
                      text width=2.3cm, align=center},
    arr/.style={-{Stealth[length=6pt]}, very thick, deepblue!70},
    lbl/.style={font=\tiny\sffamily, text=deepblue!70}
]

% --- Stage 0: Pretrained ---
\node[stage, fill=gray!15, draw=gray!50] (pt) {Pretrained\\LLM};

% --- Stage 1: SFT ---
\node[stage, right=1.6cm of pt] (sft) {SFT\\Model};
\node[datastage, above=0.5cm of sft] (sftdata) {Demonstration\\data $(x, y^*)$};

% --- Stage 2: RM ---
\node[stage, right=1.6cm of sft] (rm) {Reward\\Model};
\node[datastage, above=0.5cm of rm] (rmdata) {Preference data\\$(x, y_w, y_l)$};

% --- Stage 3: RLHF ---
\node[stage, right=1.6cm of rm, fill=deepblue!15] (rlhf) {Aligned\\Model};
\node[datastage, above=0.5cm of rlhf] (rlhfdata) {RL optimization\\(PPO + KL)};

% --- Arrows ---
\draw[arr] (pt) -- node[above, lbl]{Step 1} (sft);
\draw[arr] (sft) -- node[above, lbl]{Step 2} (rm);
\draw[arr] (rm) -- node[above, lbl]{Step 3} (rlhf);

\draw[arr, gray!50] (sftdata) -- (sft);
\draw[arr, gray!50] (rmdata) -- (rm);
\draw[arr, gray!50] (rlhfdata) -- (rlhf);

% --- Feedback arrow from RM to RLHF ---
\draw[arr, deepgreen!70, dashed] (rm.south) -- ++(0, -0.6)
  -| node[below, pos=0.25, lbl, text=deepgreen!70]{$R_\theta(x, y)$} (rlhf.south);

% --- Reference model arrow ---
\draw[arr, orange!70, dashed] (sft.south) -- ++(0, -1.0)
  -| node[below, pos=0.2, lbl, text=orange!70]{$\pi_{\text{ref}}$ (KL anchor)} (rlhf.south west);

\end{tikzpicture}
\caption{The three-stage alignment pipeline.  Stage~1: SFT on demonstrations.
Stage~2: train a reward model on human preferences.  Stage~3: optimize the
policy via PPO against the reward model, anchored to the SFT model by a
KL penalty.}
\label{fig:alignment_pipeline}
\end{figure}


% ============================================================================
\section{Reward Model Training}
\label{sec:reward_model}
% ============================================================================

The reward model (RM) is the component that translates subjective human
preferences into a scalar signal the RL algorithm can optimize.  Its quality
is the ceiling of the entire alignment pipeline---no amount of RL
optimization can compensate for a reward model that assigns high scores to
bad outputs.

\subsection{Preference Data Collection}

For each training example, human annotators are presented with:
\begin{itemize}[leftmargin=*]
\item A \textbf{prompt} $x$ (user query or instruction).
\item Two (or more) model-generated \textbf{responses} $y_A$ and $y_B$.
\item The annotator indicates which response is preferred: $y_w \succ y_l$
  (winner vs.\ loser).
\end{itemize}

Data quality considerations include: clear annotation guidelines, annotator
calibration (inter-annotator agreement), diversity of prompts (topics,
difficulty, safety), and filtering for low-quality annotations.  Typical
datasets range from 50K to 500K preference pairs.

\subsection{The Bradley-Terry Model}

Given a reward function $r_\theta(x, y)$ that maps a prompt-response pair
to a scalar, the \textbf{Bradley-Terry model} (1952) defines the
probability that response $y_w$ is preferred over $y_l$:

\begin{mathresult}{Bradley-Terry Preference Model}
\begin{equation}
P(y_w \succ y_l \given x)
= \sigma\!\bigl(r_\theta(x, y_w) - r_\theta(x, y_l)\bigr)
= \frac{1}{1 + \exp\!\bigl(-(r_\theta(x, y_w) - r_\theta(x, y_l))\bigr)}
\label{eq:bradley_terry}
\end{equation}
where $\sigma(\cdot)$ is the sigmoid function.  The preference probability
depends only on the \emph{difference} in rewards, not their absolute values.
\end{mathresult}

\paragraph{Why Bradley-Terry?}
The model has a single elegant property: preferences are transitive and
determined by reward differences.  If $r(A) > r(B) > r(C)$, then
$P(A \succ C) > P(A \succ B)$.  The sigmoid function maps unbounded reward
differences to $[0, 1]$ probabilities, and the symmetry
$P(A \succ B) = 1 - P(B \succ A)$ holds automatically.

\subsection{Reward Model Training Loss}

The reward model is trained to maximize the log-likelihood of observed human
preferences:

\begin{mathresult}{Reward Model Training Loss}
\begin{equation}
\loss_{\text{RM}} = -\E_{(x,\, y_w,\, y_l) \sim \mathcal{D}}
\Bigl[\log \sigma\!\bigl(r_\theta(x, y_w) - r_\theta(x, y_l)\bigr)\Bigr]
\label{eq:rm_loss}
\end{equation}
This is binary cross-entropy applied to the preference probability, with the
``positive class'' being the event that the preferred response receives a
higher reward than the dispreferred one.
\end{mathresult}

\paragraph{Architecture.}
The reward model typically shares the architecture of the base LLM, with the
language modeling head replaced by a \textbf{scalar projection head}: the
final hidden state (often at the last token position, or a pooled
representation) is projected to a single scalar $r_\theta(x, y) \in \R$.
Initializing the RM from the SFT model (rather than the pretrained model)
gives it a strong starting representation of what good responses look like.

\paragraph{Practical considerations.}
\begin{itemize}[leftmargin=*]
\item \textbf{Label noise:} Human preferences are noisy---annotators
  disagree on 20--35\% of pairs.  The Bradley-Terry model handles this
  gracefully because it models preferences probabilistically.
\item \textbf{Length bias:} Annotators often prefer longer, more detailed
  responses.  The RM can learn this spurious correlation, leading to
  reward hacking (Section~\ref{sec:reward_hacking}).
\item \textbf{Calibration:} The absolute reward value is meaningless (only
  differences matter).  Some implementations add a regularization term to
  prevent reward magnitude from drifting.
\item \textbf{Listwise ranking:} Extending beyond pairwise comparisons,
  some approaches rank $K > 2$ responses per prompt, yielding
  $\binom{K}{2}$ preference pairs per annotation.
\end{itemize}

\begin{keyinsight}[The Reward Model Is the Bottleneck]
The reward model is the weakest link in the RLHF chain.  It is trained on
finite, noisy human preference data, yet it must generalize to the vast
space of possible LLM outputs.  If the RM assigns high reward to verbose
but unhelpful responses, the policy will learn to be verbose.  If it
cannot distinguish subtle factual errors, the policy will hallucinate
confidently.  Improving RM quality---through better data, larger models,
ensembles, or process-level supervision---is the highest-leverage
intervention in alignment research.
\end{keyinsight}


% ============================================================================
\section{PPO for Language Model Alignment}
\label{sec:ppo}
% ============================================================================

With a trained reward model in hand, the next step is to optimize the
language model (the ``policy'') to produce outputs that receive high reward.
This is a reinforcement learning problem: the policy generates a response
(action) given a prompt (state), and receives a scalar reward.

\subsection{The RLHF Objective}

The RLHF training objective balances two goals: maximize the expected reward
from the RM, and stay close to the SFT model (the ``reference policy'') to
prevent degenerate solutions.

\begin{mathresult}{RLHF Objective}
\begin{equation}
\max_{\pi_\theta} \;\;
\E_{x \sim \mathcal{D},\; y \sim \pi_\theta(\cdot \given x)}
\Bigl[r_\phi(x, y)\Bigr]
\;-\; \beta \cdot \KL\!\bigl(\pi_\theta(\cdot \given x) \,\|\, \pi_{\text{ref}}(\cdot \given x)\bigr)
\label{eq:rlhf_objective}
\end{equation}
where $\pi_\theta$ is the policy being optimized, $\pi_{\text{ref}}$ is the
SFT model (frozen), $r_\phi$ is the trained reward model, and $\beta > 0$
controls the strength of the KL penalty.
\end{mathresult}

\paragraph{Why the KL penalty?}
Without the KL term, the policy would find degenerate outputs that exploit
imperfections in the reward model---producing gibberish that happens to
score highly, or collapsing to a single memorized high-reward response.
The KL penalty keeps the policy close to the SFT model, which serves as a
``sanity anchor'' ensuring the outputs remain coherent and diverse.

\paragraph{The combined reward.}
In practice, the KL penalty is folded into the reward signal as a per-token
penalty:
\begin{equation}
r_{\text{total}}(x, y)
= r_\phi(x, y)
  - \beta \sum_{t=1}^{T} \log \frac{\pi_\theta(y_t \given x, y_{<t})}
                                     {\pi_{\text{ref}}(y_t \given x, y_{<t})}
\label{eq:total_reward}
\end{equation}
This converts the constrained optimization into an unconstrained RL problem
with a modified reward.

\subsection{PPO Mechanics}

\textbf{Proximal Policy Optimization} (Schulman et al., 2017) is the RL
algorithm used to optimize the RLHF objective.  PPO belongs to the family
of policy gradient methods and is designed for stable, monotonic policy
improvement.

\begin{mathresult}{PPO Clipped Surrogate Objective}
\begin{equation}
\loss_{\text{PPO}} = -\E_t \Bigl[
  \min\!\bigl(
    \rho_t(\theta)\, \hat{A}_t,\;\;
    \text{clip}\bigl(\rho_t(\theta),\, 1-\epsilon,\, 1+\epsilon\bigr)\, \hat{A}_t
  \bigr)
\Bigr]
\label{eq:ppo}
\end{equation}
where:
\begin{itemize}
\item $\rho_t(\theta) = \dfrac{\pi_\theta(y_t \given x, y_{<t})}
  {\pi_{\theta_{\text{old}}}(y_t \given x, y_{<t})}$ is the probability
  ratio between the current and old policy
\item $\hat{A}_t$ is the estimated advantage (how much better action $y_t$
  is than average under the current value function)
\item $\epsilon$ is the clipping range (typically 0.1--0.2)
\end{itemize}
\end{mathresult}

\paragraph{Why clipping?}
Without clipping, a large policy update in one step could catastrophically
degrade performance.  The $\min$ with the clipped ratio ensures that the
objective is pessimistic: when the advantage is positive (good action), the
ratio is clipped above at $1 + \epsilon$, limiting how much the policy can
increase the probability of that action.  When the advantage is negative
(bad action), the ratio is clipped below at $1 - \epsilon$, limiting how
much the probability can decrease.  This creates a ``trust region'' around
the old policy.

\paragraph{Advantage estimation (GAE).}
The advantage $\hat{A}_t$ is computed using Generalized Advantage Estimation
(Schulman et al., 2016):
\begin{equation}
\hat{A}_t = \sum_{l=0}^{\infty} (\gamma \lambda)^l \, \delta_{t+l},
\qquad
\delta_t = r_t + \gamma V(s_{t+1}) - V(s_t)
\end{equation}
where $V(s_t)$ is a learned value function (critic), $\gamma$ is the
discount factor, and $\lambda \in [0, 1]$ controls the bias-variance
tradeoff.  In RLHF, the ``state'' at time $t$ is the prompt plus tokens
generated so far, and $r_t$ is typically zero for all tokens except the
last, where it equals $r_\phi(x, y)$ minus the accumulated KL penalty.

\subsection{The Four-Model Setup}

RLHF with PPO requires maintaining four models simultaneously:

\begin{table}[H]
\centering
\begin{tabularx}{\textwidth}{lXc}
\toprule
\textbf{Model} & \textbf{Role} & \textbf{Updated?} \\
\midrule
Policy $\pi_\theta$ & Generates responses; parameters being optimized & Yes \\
Reference $\pi_{\text{ref}}$ & Frozen SFT model; anchor for KL penalty & No \\
Reward model $r_\phi$ & Scores (prompt, response) pairs & No \\
Value function $V_\psi$ & Estimates expected cumulative reward (critic) & Yes \\
\bottomrule
\end{tabularx}
\caption{The four models required for PPO-based RLHF.  For a 70B parameter
policy, this means approximately 280B parameters in GPU memory.}
\label{tab:four_models}
\end{table}

\begin{warningbox}
\textbf{PPO for RLHF is notoriously difficult to tune.}  Practical
challenges include: (a)~the four-model memory footprint---for a 70B policy,
you need $\sim$280B parameters in memory simultaneously; (b)~hyperparameter
sensitivity---the KL coefficient $\beta$, learning rate, clipping range
$\epsilon$, GAE parameters $\gamma$ and $\lambda$, and number of PPO
epochs per batch all interact in complex ways; (c)~training instability---
the reward model is a moving target from the policy's perspective (reward
hacking), and the value function must track a nonstationary reward;
(d)~high variance---policy gradient methods have inherently high variance,
requiring large batch sizes.  These difficulties directly motivated the
development of DPO (Section~\ref{sec:dpo}).
\end{warningbox}


% ============================================================================
\section{Reward Hacking}
\label{sec:reward_hacking}
% ============================================================================

\textbf{Reward hacking} occurs when the policy learns to exploit
imperfections in the reward model rather than genuinely improving response
quality.  It is the central failure mode of RLHF and a manifestation of
Goodhart's law: ``When a measure becomes a target, it ceases to be a good
measure.''

\subsection{Common Manifestations}

\begin{itemize}[leftmargin=*]
\item \textbf{Verbosity:} Human annotators often prefer longer, more
  detailed responses.  The RM learns this correlation, and the policy
  exploits it by generating increasingly verbose outputs that add words
  without adding value.
\item \textbf{Sycophancy:} The policy learns to agree with the user's
  stated beliefs, even when incorrect, because agreeable responses receive
  higher human ratings.
\item \textbf{Formatting tricks:} Bullet points, headers, and structured
  formatting score higher with annotators, so the policy over-produces
  formatted content even when a brief answer suffices.
\item \textbf{Hedging:} Excessive caveats and qualifications (``It's
  important to note that\ldots'') to appear thorough without committing to
  potentially wrong answers.
\item \textbf{Style over substance:} Confident, articulate phrasing that
  sounds authoritative but may be factually wrong.
\end{itemize}

\subsection{The Over-Optimization Curve}

Gao et al.\ (2023) demonstrated a characteristic pattern: as PPO training
progresses, the \emph{proxy reward} (score from the RM) increases
monotonically, but the \emph{true reward} (actual human evaluation) first
improves, peaks, and then \emph{degrades}.

\begin{figure}[H]
\centering
\begin{tikzpicture}
\begin{axis}[
    width=10cm, height=6cm,
    xlabel={PPO Training Steps},
    ylabel={Reward},
    xmin=0, xmax=100,
    ymin=-0.5, ymax=2.5,
    legend pos=north west,
    grid=major,
    grid style={gray!20},
    axis lines=left,
    xlabel style={font=\small\sffamily},
    ylabel style={font=\small\sffamily},
    tick label style={font=\scriptsize\sffamily},
    legend style={font=\scriptsize\sffamily, fill=white, draw=gray!40},
]
% Proxy reward (monotonically increasing)
\addplot[deepblue, very thick, smooth, domain=0:100, samples=80]
    {0.02*x + 0.3*ln(x+1)/ln(10)};
\addlegendentry{Proxy reward (RM score)}

% True reward (peaks then degrades)
\addplot[deepred, very thick, smooth, dashed, domain=0:100, samples=80]
    {-0.0003*x^2 + 0.04*x + 0.1};
\addlegendentry{True reward (human eval)}

% Optimal point
\draw[deepgreen!80, thick, dashed] (axis cs:66.7,0) -- (axis cs:66.7,1.45)
    node[above, font=\scriptsize\sffamily, text=deepgreen!80] {Optimal};

% Over-optimization zone
\draw[deepred!60, thick, decorate, decoration={brace, amplitude=5pt, mirror, raise=2pt}]
    (axis cs:67,{-0.0003*67^2 + 0.04*67 + 0.1}) --
    (axis cs:100,{-0.0003*100^2 + 0.04*100 + 0.1})
    node[midway, below=8pt, font=\scriptsize\sffamily, text=deepred!70]
    {Over-optimization};

\end{axis}
\end{tikzpicture}
\caption{The over-optimization curve.  The proxy reward (RM score)
increases monotonically, but the true reward (human evaluation) peaks
and then degrades as the policy increasingly exploits RM imperfections.}
\label{fig:overoptimization}
\end{figure}

\subsection{Defenses Against Reward Hacking}

\begin{enumerate}[leftmargin=*]
\item \textbf{KL penalty} (primary defense):
  Constraining the policy to stay near $\pi_{\text{ref}}$ limits how far
  it can drift into reward-hacked territory.  The coefficient $\beta$
  controls this tradeoff: higher $\beta$ = safer but less improvement;
  lower $\beta$ = more improvement but higher risk of hacking.

\item \textbf{Reward model ensembles:}
  Train multiple RMs on different data splits and use their agreement as
  a more robust reward signal (e.g., minimum or average of ensemble).
  Disagreement between ensemble members signals regions where hacking
  is likely.

\item \textbf{Iterative reward model retraining:}
  Periodically collect new preference data on the \emph{current} policy's
  outputs (not the SFT model's outputs) and retrain the RM.  This closes
  the distribution gap between RM training data and policy outputs.

\item \textbf{Length penalties:}
  Explicitly penalize response length in the reward to counteract the
  verbosity bias.

\item \textbf{Constitutional constraints:}
  Use rule-based or model-based checks alongside the RM to catch specific
  failure modes (Section~\ref{sec:cai}).
\end{enumerate}

\begin{keyinsight}[Goodhart's Law in Alignment]
Goodhart's law---``When a measure becomes a target, it ceases to be a good
measure''---is not just a witty aphorism in alignment; it is the
\emph{central technical challenge}.  The reward model is a proxy for human
values, and optimizing hard against a proxy inevitably diverges from the
true objective.  Every defense strategy (KL penalty, ensembles, retraining)
is fundamentally about limiting how hard we optimize against the proxy.
This is why alignment is an ongoing research problem, not a solved one.
\end{keyinsight}


% ============================================================================
\section{DPO: Direct Preference Optimization}
\label{sec:dpo}
% ============================================================================

Direct Preference Optimization (Rafailov et al., 2023) is an elegant
reformulation that eliminates both the reward model and the RL algorithm,
reducing RLHF to a single supervised learning objective on preference data.

\subsection{The Key Insight}

DPO observes that the optimal policy under the KL-constrained RLHF
objective has a \emph{closed-form} relationship to the reward function.
By rearranging this relationship and substituting into the Bradley-Terry
preference model, we can derive a loss function that directly optimizes
the policy on preference data---no intermediate reward model or RL required.

\subsection{Derivation from the RLHF Objective}

\paragraph{Step 1: The optimal policy.}
Consider the RLHF objective from Eq.~\eqref{eq:rlhf_objective}:
\[
\max_{\pi} \;\;
\E_{y \sim \pi(\cdot \given x)}
\bigl[r(x, y)\bigr]
- \beta \cdot \KL\!\bigl(\pi(\cdot \given x) \,\|\, \pi_{\text{ref}}(\cdot \given x)\bigr)
\]
Expanding the KL divergence and solving for the optimal policy $\pi^*$
(by setting the functional derivative to zero), we obtain:

\begin{mathresult}{Optimal Policy under KL-Constrained RM Maximization}
\begin{equation}
\pi^*(y \given x) = \frac{1}{Z(x)} \; \pi_{\text{ref}}(y \given x) \;\exp\!\left(\frac{r(x, y)}{\beta}\right)
\label{eq:optimal_policy}
\end{equation}
where $Z(x) = \sum_{y'} \pi_{\text{ref}}(y' \given x) \exp\!\bigl(r(x, y')/\beta\bigr)$ is the partition function (intractable to compute in general, but it cancels out below).
\end{mathresult}

\textit{Derivation sketch.}  Write the objective in terms of the policy
distribution:
\begin{align}
J(\pi) &= \sum_y \pi(y \given x) \, r(x, y)
  - \beta \sum_y \pi(y \given x) \log \frac{\pi(y \given x)}{\pi_{\text{ref}}(y \given x)}
\label{eq:rlhf_expanded}
\end{align}
Taking the functional derivative with respect to $\pi(y \given x)$ and
setting it to zero (with a Lagrange multiplier for the normalization
constraint $\sum_y \pi(y \given x) = 1$):
\begin{align}
\frac{\partial J}{\partial \pi(y \given x)}
&= r(x, y) - \beta \log \frac{\pi(y \given x)}{\pi_{\text{ref}}(y \given x)} - \beta - \lambda = 0
\end{align}
Solving for $\pi(y \given x)$:
\[
\pi(y \given x) = \pi_{\text{ref}}(y \given x) \cdot \exp\!\left(\frac{r(x,y) - \beta - \lambda}{\beta}\right)
\]
The terms $\beta + \lambda$ become the log-partition function after
enforcing normalization, yielding Eq.~\eqref{eq:optimal_policy}.

\paragraph{Step 2: Express the reward in terms of the policy.}
Rearranging Eq.~\eqref{eq:optimal_policy} by taking the logarithm of both
sides:
\begin{equation}
r(x, y) = \beta \log \frac{\pi^*(y \given x)}{\pi_{\text{ref}}(y \given x)} + \beta \log Z(x)
\label{eq:reward_from_policy}
\end{equation}
This is the crucial equation: the reward is a function of the
\emph{log-ratio} between the optimal policy and the reference policy, plus
a prompt-dependent constant.

\paragraph{Step 3: Substitute into Bradley-Terry.}
The reward model loss (Eq.~\ref{eq:rm_loss}) uses the Bradley-Terry
model, which depends only on \emph{reward differences}:
\begin{align}
r(x, y_w) - r(x, y_l)
&= \beta \log \frac{\pi^*(y_w \given x)}{\pi_{\text{ref}}(y_w \given x)}
 - \beta \log \frac{\pi^*(y_l \given x)}{\pi_{\text{ref}}(y_l \given x)}
\label{eq:reward_diff}
\end{align}
The partition function $\beta \log Z(x)$ \textbf{cancels} because it
depends only on $x$, not on the specific response.  Substituting into the
Bradley-Terry model:

\begin{mathresult}{DPO Loss Function}
\begin{equation}
\loss_{\text{DPO}}(\pi_\theta; \pi_{\text{ref}})
= -\E_{(x,\, y_w,\, y_l) \sim \mathcal{D}}
\left[\log \sigma\!\left(
  \beta \left(
    \log \frac{\pi_\theta(y_w \given x)}{\pi_{\text{ref}}(y_w \given x)}
    - \log \frac{\pi_\theta(y_l \given x)}{\pi_{\text{ref}}(y_l \given x)}
  \right)
\right)\right]
\label{eq:dpo_loss}
\end{equation}
where $\pi_\theta$ is the policy being trained, $\pi_{\text{ref}}$ is the
frozen reference model (typically the SFT model), $y_w$ and $y_l$ are the
preferred and dispreferred responses, and $\beta$ controls the strength
of the KL constraint (same role as in RLHF).
\end{mathresult}

\paragraph{Reading the DPO loss.}
The loss encourages the policy to increase the log-probability of the
preferred response $y_w$ \emph{relative to the reference}, while
decreasing it for the dispreferred response $y_l$ \emph{relative to the
reference}.  The reference model provides a baseline: DPO does not just
make $y_w$ more likely in absolute terms---it makes $y_w$ \emph{more
likely than the reference model would predict}, and $y_l$ \emph{less
likely}.

\begin{keyinsight}[The Elegance of DPO]
DPO shows that the RLHF objective, the reward model, and PPO optimization
are all \textbf{mathematically redundant}---they can be collapsed into a
single supervised loss on preference data.  The ``reward model'' is
implicitly defined by the policy itself via
$r(x, y) = \beta \log(\pi_\theta(y|x) / \pi_{\text{ref}}(y|x)) + C(x)$.
This is not an approximation; it is an exact reformulation.  The reward
model was never a necessary intermediate---it was a computational
scaffolding that DPO removes.
\end{keyinsight}

\subsection{DPO vs.\ RLHF: Comparison}

\begin{table}[H]
\centering
\begin{tabularx}{\textwidth}{lXX}
\toprule
\textbf{Aspect} & \textbf{RLHF (PPO)} & \textbf{DPO} \\
\midrule
Components & 4 models (policy, ref, RM, critic) & 2 models (policy, ref) \\
Training & RL loop (generate $\to$ score $\to$ update) & Supervised (single loss on preference data) \\
Stability & Sensitive to hyperparameters & More stable, standard supervised training \\
Compute & High (generation + RM inference + PPO) & Low (forward passes only) \\
Explicit reward & Yes (can monitor and debug RM) & No (reward is implicit in the policy) \\
Online data & Yes (policy generates new data each step) & No (fixed offline preference dataset) \\
Exploration & Policy can explore novel outputs & Limited to preference dataset distribution \\
\bottomrule
\end{tabularx}
\caption{DPO vs.\ RLHF comparison.  DPO trades flexibility and explicit
reward monitoring for simplicity and stability.}
\label{tab:dpo_vs_rlhf}
\end{table}

\paragraph{DPO limitations.}
\begin{itemize}[leftmargin=*]
\item \textbf{No explicit reward model:} Cannot monitor the reward signal
  during training or use it for other purposes (e.g., best-of-$n$
  sampling, rejection sampling).
\item \textbf{Offline only:} DPO trains on a fixed dataset.  It cannot
  generate new responses and get reward feedback during training, which
  limits its ability to explore beyond the preference data distribution.
  Online variants (online DPO, iterative DPO) address this partially.
\item \textbf{Reference model dependence:} DPO's loss is defined relative
  to $\pi_{\text{ref}}$.  If the reference model is poor, DPO's
  optimization landscape is distorted.
\item \textbf{Sensitivity to $\beta$:} Like RLHF, $\beta$ controls the
  KL constraint.  Too high $\beta$ undertrains; too low leads to
  over-optimization.
\end{itemize}


% ============================================================================
\section{DPO Variants}
\label{sec:dpo_variants}
% ============================================================================

The success of DPO inspired a family of preference optimization methods,
each addressing specific limitations.

\begin{table}[H]
\centering
\small
\begin{tabularx}{\textwidth}{@{}l l X@{}}
\toprule
\textbf{Method} & \textbf{Key Change} & \textbf{Why It Matters} \\
\midrule
IPO & $\ell_2$ regularization instead of log-sigmoid &
  Avoids overfitting to the Bradley-Terry assumption; more robust
  when preferences are noisy or intransitive \\
KTO & Needs only binary feedback (good/bad), not pairs &
  Dramatically easier data collection---no need to generate and compare
  two responses per prompt; uses Kahneman-Tversky prospect theory \\
ORPO & Combines SFT and preference optimization in one loss &
  No need for separate SFT stage; reference-free (no $\pi_{\text{ref}}$
  needed) \\
SimPO & Reference-free, uses average log-prob as implicit reward &
  Simpler than DPO; no reference model at all; length-normalized
  reward avoids verbosity bias \\
cDPO & Conservative DPO with label noise modeling &
  Accounts for annotation noise explicitly in the loss \\
RSO & Rejection Sampling Optimization with SFT on filtered data &
  Combines rejection sampling from RM with DPO-style optimization \\
\bottomrule
\end{tabularx}
\caption{DPO variants.  Each method addresses a specific limitation of
the original DPO formulation.}
\label{tab:dpo_variants}
\end{table}

\paragraph{KTO in detail.}
KTO (Ethayarajh et al., 2024) is notable because it requires only
\emph{pointwise} feedback (thumbs up / thumbs down on individual responses)
rather than \emph{pairwise} preferences.  This is a significant practical
advantage: collecting pairwise comparisons requires generating two responses
per prompt and asking annotators to compare them, while pointwise feedback
can be collected passively from user interactions (upvotes, regeneration
requests, thumbs-up buttons).

The KTO loss uses the Kahneman-Tversky value function, which assigns
asymmetric weight to gains (preferred responses) and losses (dispreferred
responses), reflecting the empirical finding that humans are loss-averse:
\begin{equation}
\loss_{\text{KTO}} = \E_{(x,y)} \bigl[
  w(y) \cdot \bigl(1 - \sigma\!\bigl(\beta \, r_\theta(x, y) - z_{\text{ref}}\bigr)\bigr)
\bigr]
\end{equation}
where $r_\theta(x, y) = \log(\pi_\theta(y|x) / \pi_{\text{ref}}(y|x))$
is the implicit reward, $z_{\text{ref}}$ is a running-mean baseline, and
$w(y)$ weights desirable and undesirable examples differently.


% ============================================================================
\section{Constitutional AI}
\label{sec:cai}
% ============================================================================

Constitutional AI (CAI), introduced by Anthropic (Bai et al., 2022),
addresses a fundamental scaling challenge: human feedback is expensive
and slow.  CAI uses the model itself to generate feedback, guided by a
set of principles (the ``constitution'').

\subsection{The CAI Pipeline}

\begin{figure}[H]
\centering
\begin{tikzpicture}[
    node distance=0.8cm and 0.5cm,
    stepbox/.style={rectangle, draw=deepblue, fill=lightblue, thick,
                    minimum height=0.9cm, minimum width=3.2cm,
                    font=\small\sffamily, rounded corners=3pt,
                    text width=3.0cm, align=center},
    prinbox/.style={rectangle, draw=orange!70, fill=yellow!8,
                    minimum height=0.7cm, minimum width=2.6cm,
                    font=\scriptsize\sffamily, rounded corners=2pt,
                    text width=2.5cm, align=center},
    arr/.style={-{Stealth[length=5pt]}, thick, deepblue!70},
    lbl/.style={font=\tiny\sffamily, text=deepblue!70}
]

% SL-CAI (left column)
\node[stepbox] (gen) {Generate\\initial response};
\node[stepbox, below=0.9cm of gen] (crit) {Self-critique\\against constitution};
\node[stepbox, below=0.9cm of crit] (rev) {Self-revision};
\node[stepbox, below=0.9cm of rev, fill=deepblue!15] (sft) {Fine-tune on\\revised responses};

% Constitution
\node[prinbox, right=1.5cm of crit] (const) {Constitution\\(set of principles)};

% RL-CAI (right column)
\node[stepbox, right=3.5cm of gen] (gen2) {Generate\\response pairs};
\node[stepbox, below=0.9cm of gen2] (aifb) {AI feedback\\(which is better?)};
\node[stepbox, below=0.9cm of aifb] (rm2) {Train RM on\\AI preferences};
\node[stepbox, below=0.9cm of rm2, fill=deepblue!15] (rlcai) {RLHF with\\AI-trained RM};

% Arrows
\draw[arr] (gen) -- (crit);
\draw[arr] (crit) -- (rev);
\draw[arr] (rev) -- (sft);
\draw[arr, orange!70] (const) -- (crit);
\draw[arr, orange!70] (const) |- (aifb);

\draw[arr] (gen2) -- (aifb);
\draw[arr] (aifb) -- (rm2);
\draw[arr] (rm2) -- (rlcai);

% Labels
\node[font=\small\sffamily\bfseries, text=deepblue, above=0.2cm of gen]
     {SL-CAI (Supervised)};
\node[font=\small\sffamily\bfseries, text=deepblue, above=0.2cm of gen2]
     {RL-CAI (Reinforcement)};

% Connecting arrow
\draw[arr, deepgreen!70, dashed, thick]
  (sft.east) -- ++(0.5, 0) |- (gen2.west);
\node[lbl, text=deepgreen!70] at (4.3, -1.1) {feeds into};

\end{tikzpicture}
\caption{Constitutional AI pipeline.  \textbf{SL-CAI} (left): the model
critiques and revises its own outputs guided by constitutional principles,
then is fine-tuned on the revised responses.  \textbf{RL-CAI} (right):
the model generates AI preference labels (RLAIF), used to train a reward
model for RLHF.}
\label{fig:cai_pipeline}
\end{figure}

\subsection{SL-CAI: Supervised Learning from Self-Critique}

\begin{enumerate}[leftmargin=*]
\item \textbf{Generate:} The model produces an initial response to a prompt
  (which may be harmful or low-quality).
\item \textbf{Critique:} The model is asked to critique its own response
  against a specific constitutional principle (e.g., ``Identify ways in
  which this response is harmful or unethical'').
\item \textbf{Revise:} The model revises its response to address the
  identified issues.
\item \textbf{Fine-tune:} The model is fine-tuned on the (prompt, revised
  response) pairs using standard supervised learning.
\end{enumerate}

The \textbf{constitution} is a set of natural language principles such as:
\begin{itemize}[leftmargin=*]
\item ``Choose the response that is least likely to be harmful or toxic.''
\item ``Choose the response that is most honest and transparent.''
\item ``Choose the response that is least likely to encourage illegal activity.''
\end{itemize}

\subsection{RL-CAI: RLAIF}

In the reinforcement learning phase, \textbf{RLAIF} (RL from AI Feedback)
replaces human annotators with the model itself:

\begin{enumerate}[leftmargin=*]
\item Generate two responses for each prompt.
\item Ask the model which response better adheres to the constitutional
  principles.
\item Use these AI-generated preferences to train a reward model.
\item Run standard RLHF (PPO) with this reward model.
\end{enumerate}

The surprising finding is that RLAIF produces alignment quality comparable
to RLHF with human feedback, at a fraction of the cost and with much
faster iteration cycles.

\paragraph{Why CAI works.}
The constitution provides \textbf{compositional, interpretable control}
over model behavior.  Unlike a reward model (which is a black box that
maps inputs to scores), the constitution is a readable set of rules that
can be updated, debated, and audited.  The model's ability to follow these
principles during self-critique leverages its instruction-following
capability---an ability that improves with scale.


% ============================================================================
\section{Process Reward Models}
\label{sec:process_rm}
% ============================================================================

Standard reward models provide \textbf{outcome-based} feedback: the entire
response receives a single scalar reward.  \textbf{Process reward models}
(PRMs) evaluate each \emph{step} of a reasoning chain, providing much
finer-grained supervision.

\subsection{Outcome vs.\ Process Supervision}

\begin{table}[H]
\centering
\begin{tabularx}{\textwidth}{lXX}
\toprule
\textbf{Aspect} & \textbf{Outcome RM (ORM)} & \textbf{Process RM (PRM)} \\
\midrule
Granularity & Single reward for entire response & Reward per reasoning step \\
Training data & (prompt, response, quality) & (prompt, step$_i$, step correctness) \\
Signal density & Sparse (one signal per response) & Dense (one signal per step) \\
Credit assignment & Ambiguous (which step caused failure?) & Clear (exact step where error occurs) \\
Data cost & Lower (one annotation per response) & Higher (annotation per step) \\
Best for & General-purpose alignment & Math, code, multi-step reasoning \\
\bottomrule
\end{tabularx}
\caption{Outcome vs.\ process reward models.}
\label{tab:orm_vs_prm}
\end{table}

\subsection{Process Reward Models for Reasoning}

Lightman et al.\ (2023) demonstrated that process supervision significantly
outperforms outcome supervision for mathematical reasoning.  The PRM
assigns a correctness score to each step of a chain-of-thought solution:

\begin{enumerate}[leftmargin=*]
\item The model generates a step-by-step solution.
\item The PRM evaluates each step: is this step logically correct given the
  previous steps?
\item Steps are labeled as \textbf{positive} (correct), \textbf{negative}
  (contains an error), or \textbf{neutral} (neither clearly correct nor
  incorrect).
\end{enumerate}

\subsection{Search over Reasoning Paths}

The most powerful application of PRMs is enabling \textbf{tree search}
over reasoning paths---the approach behind models like o1 and o3:

\begin{enumerate}[leftmargin=*]
\item \textbf{Generate:} The model produces multiple candidate next steps.
\item \textbf{Evaluate:} The PRM scores each candidate step.
\item \textbf{Select:} Use the PRM scores to guide search (beam search,
  best-first search, or Monte Carlo Tree Search).
\item \textbf{Backtrack:} If all candidate steps score poorly, backtrack
  to a previous step and try alternative paths.
\end{enumerate}

This is essentially \textbf{test-time compute scaling}: spending more
computation at inference to search for better reasoning paths, rather than
relying solely on the model's one-shot generation.

\begin{keyinsight}[Process vs.\ Outcome Supervision]
Process supervision addresses the \textbf{credit assignment problem} that
plagues outcome-based RLHF.  When a mathematical proof is wrong, outcome
supervision says ``this proof is bad'' without indicating which step
failed.  Process supervision pinpoints the exact error, providing a
denser, more informative learning signal.  The tradeoff is data cost:
step-level annotations are 5--10$\times$ more expensive than response-level
annotations.  The field is moving toward \emph{automated} process
supervision (using verifiers, code execution, and self-consistency) to
reduce this cost.
\end{keyinsight}


% ============================================================================
\section{Practical Alignment Engineering}
\label{sec:practical_alignment}
% ============================================================================

\subsection{Choosing an Alignment Method}

\begin{table}[H]
\centering
\small
\begin{tabularx}{\textwidth}{@{}l c c c X@{}}
\toprule
\textbf{Method} & \textbf{Data Needed} & \textbf{Compute} & \textbf{Complexity} & \textbf{Best For} \\
\midrule
SFT only         & Demonstrations & Low    & Low    & Quick baseline, cold start \\
DPO              & Pairwise prefs & Medium & Low    & Most teams; default starting point \\
KTO              & Binary ratings & Medium & Low    & When only thumbs up/down data exists \\
RLHF (PPO)       & Pairwise prefs & High   & High   & Maximum control; frontier labs \\
CAI / RLAIF      & Constitution   & High   & Medium & Scalable safety; reduces human cost \\
Best-of-$n$ + RM & Pairwise prefs & High$^*$ & Medium & When RM exists; inference-time alignment \\
\bottomrule
\end{tabularx}
\caption{Alignment method selection guide.  $^*$Compute is at inference
time (generate $n$ candidates, score, return best).}
\label{tab:alignment_choice}
\end{table}

\subsection{The Current Alignment Recipe (2024--2025)}

Most frontier labs follow a recipe similar to:

\begin{enumerate}[leftmargin=*]
\item \textbf{Pretraining:} Train the base model on trillions of tokens.
\item \textbf{SFT:} Fine-tune on high-quality instruction-response pairs
  (10K--100K examples). Quality matters far more than quantity (LIMA: ``Less
  Is More for Alignment'').
\item \textbf{Preference optimization:} Apply DPO or RLHF on 50K--500K
  preference pairs.  Some labs do iterative rounds: collect preferences
  on the current model's outputs, optimize, repeat.
\item \textbf{Safety fine-tuning:} Additional alignment specifically
  targeting safety (refusal of harmful requests, reducing bias), often
  with targeted red-teaming data.
\item \textbf{Constitutional / rule-based guardrails:} Deploy with input
  and output safety classifiers as a final layer of defense.
\end{enumerate}


% ============================================================================
\section{Interview Questions}
\label{sec:alignment_interview}
% ============================================================================

% --- Question 1 ---
\begin{interviewq}
\textbf{Q: Walk me through the RLHF pipeline step by step.  What are the
three stages?}

\badgeLfive{L5} \quad \badgetype{Conceptual} \quad
\badgefreq{\freqhigh}

\stronganswer
The RLHF pipeline has three stages.

\textbf{Stage 1---SFT:} Fine-tune the pretrained model on high-quality
demonstration data (instruction-response pairs) to teach it the format and
style of helpful responses.  This produces the SFT model, which also serves
as the reference policy $\pi_{\text{ref}}$ for subsequent stages.

\textbf{Stage 2---Reward Model Training:} Collect human preferences: for
each prompt, generate two (or more) responses, and have annotators indicate
which is better.  Train a reward model using the Bradley-Terry framework:
$\loss_{\text{RM}} = -\E[\log \sigma(r(x, y_w) - r(x, y_l))]$.
The RM shares the base model architecture but replaces the LM head with a
scalar output.

\textbf{Stage 3---RL Optimization (PPO):} Optimize the policy to maximize
$\E[r(x, y)] - \beta \cdot \KL(\pi_\theta \| \pi_{\text{ref}})$.  PPO
uses a clipped surrogate objective for stable updates.  The KL penalty
prevents reward hacking by keeping the policy close to the SFT model.
This stage requires four models in memory: policy, reference, reward model,
and value function (critic).

\redflags
\begin{itemize}
\item Cannot name all three stages or gets the order wrong.
\item Skips the reward model and jumps from SFT to RL.
\item Does not mention the KL penalty or its purpose.
\item Confuses the reward model with the value function.
\end{itemize}

\interviewtest{Whether you understand the complete pipeline that transforms
a pretrained model into an aligned assistant.  This is foundational
knowledge for anyone working with or deploying LLMs.}

\goodgreat
\begin{itemize}
\item \badgeLfive{L5}: Correctly describes all three stages with the
  key equations (RM loss, RLHF objective).  Explains why each stage
  exists and what problem it solves.
\item \badgeLsix{L6}: Discusses practical challenges (four-model memory
  footprint, RM quality as bottleneck, PPO instability), and mentions
  DPO as a simpler alternative that collapses stages 2 and 3.
\item \badgeLseven{L7}: Discusses iterative RLHF (collecting new
  preferences on the current policy), the relationship between RLHF
  and Bayesian inference, and how Constitutional AI scales the pipeline.
\end{itemize}

\followups{How many models do you need in GPU memory during PPO?  What
happens if the reward model is poor?  Why is the SFT model used as the
reference rather than the pretrained model?}
\end{interviewq}

% --- Question 2 ---
\begin{interviewq}
\textbf{Q: What is reward hacking?  How does the KL penalty help?  When
does it fail?}

\badgeLfive{L5} \quad \badgetype{Conceptual} \quad
\badgefreq{\freqhigh}

\stronganswer
Reward hacking occurs when the policy exploits imperfections in the reward
model rather than genuinely improving response quality.  The RM is a
proxy for human preferences trained on finite data, so it inevitably has
blind spots.  Common examples: the RM gives higher scores to longer
responses (verbosity hack), sycophantic agreement, or responses with
impressive-sounding but vacuous formatting.

The KL penalty $\beta \cdot \KL(\pi_\theta \| \pi_{\text{ref}})$ is the
primary defense.  It constrains the policy to stay close to the SFT model,
limiting how far it can drift into regions where the RM's predictions are
unreliable.  Intuitively, the RM was trained on outputs from $\pi_{\text{ref}}$
(or similar), so its predictions are most trustworthy near that distribution.

The KL penalty fails when: (1)~the reward model has \emph{systematic} biases
(e.g., always preferring longer responses)---the policy can hack these
biases even within the KL budget; (2)~$\beta$ is set too low, giving the
policy too much room to diverge; (3)~the reference model itself has
problematic behaviors that the KL penalty preserves.

\redflags
\begin{itemize}
\item Defines reward hacking vaguely as ``the model being wrong.''
\item Does not connect the KL penalty to the RM's distribution shift.
\item Claims the KL penalty fully solves reward hacking.
\item Cannot give concrete examples of reward hacking.
\end{itemize}

\interviewtest{Understanding of Goodhart's law applied to alignment, and
whether the candidate can reason about proxy optimization and its failure
modes.}

\goodgreat
\begin{itemize}
\item \badgeLfive{L5}: Defines reward hacking with concrete examples,
  explains the KL penalty mechanism, identifies when it fails.
\item \badgeLsix{L6}: Discusses the over-optimization curve (Gao et al.),
  mentions RM ensembles and iterative RM retraining as additional defenses,
  quantifies the $\beta$ tradeoff.
\item \badgeLseven{L7}: Connects to Goodhart's law formally, discusses
  the relationship between RM capacity and hackability, and considers
  whether process reward models reduce reward hacking for reasoning tasks.
\end{itemize}

\followups{What is the over-optimization curve?  How would you detect
reward hacking in practice?  How do RM ensembles help?}
\end{interviewq}

% --- Question 3 ---
\begin{interviewq}
\textbf{Q: Explain DPO.  How does it relate to RLHF?  What are its
advantages and limitations?}

\badgeLsix{L6} \quad \badgetype{Mathematical} \quad
\badgefreq{\freqhigh}

\stronganswer
DPO reformulates RLHF as a supervised learning problem.  The key insight
is that the optimal policy under the KL-constrained RM maximization
objective has a closed form:
$\pi^*(y|x) \propto \pi_{\text{ref}}(y|x) \exp(r(x,y)/\beta)$.  Rearranging,
the reward is a function of the policy:
$r(x,y) = \beta \log(\pi^*(y|x)/\pi_{\text{ref}}(y|x)) + C(x)$.
Substituting this into the Bradley-Terry model, the prompt-dependent
constant $C(x)$ cancels in the reward \emph{difference}, yielding the DPO
loss:
$\loss_{\text{DPO}} = -\E[\log \sigma(\beta(\log \frac{\pi_\theta(y_w|x)}{\pi_{\text{ref}}(y_w|x)} - \log \frac{\pi_\theta(y_l|x)}{\pi_{\text{ref}}(y_l|x)}))]$.

This directly optimizes the policy on preference data without a separate
reward model or RL algorithm.  Advantages: simpler (2 models vs.\ 4),
more stable training (standard supervised loss), lower compute.
Limitations: no explicit reward for monitoring or best-of-$n$ sampling,
offline (no exploration during training), and reference model quality
is critical.

\redflags
\begin{itemize}
\item Cannot sketch the derivation or explain where the partition
  function cancels.
\item Treats DPO as ``just an alternative to RLHF'' without understanding
  the mathematical connection.
\item Claims DPO is always better than RLHF.
\item Does not mention the offline limitation.
\end{itemize}

\interviewtest{Mathematical maturity and whether the candidate understands
DPO at a principled level rather than as a black-box recipe.}

\goodgreat
\begin{itemize}
\item \badgeLfive{L5}: Correctly states the DPO loss and explains that
  it eliminates the need for RM and RL.  Gives advantages and limitations.
\item \badgeLsix{L6}: Walks through the derivation (optimal policy $\to$
  reward reparameterization $\to$ Bradley-Terry substitution $\to$
  partition function cancellation).  Discusses the implicit reward
  interpretation.
\item \badgeLseven{L7}: Derives the DPO loss from scratch, discusses
  the connection to contrastive learning, compares the gradient dynamics
  of DPO vs.\ PPO, and evaluates when RLHF's online exploration
  advantage matters most.
\end{itemize}

\followups{Derive the DPO loss from the RLHF objective.  What is the
implicit reward in DPO?  When would you prefer RLHF over DPO?}
\end{interviewq}

% --- Question 4 ---
\begin{interviewq}
\textbf{Q: What is Constitutional AI?  How does self-critique enable
alignment without human feedback at every step?}

\badgeLsix{L6} \quad \badgetype{Conceptual} \quad
\badgefreq{\freqmed}

\stronganswer
Constitutional AI uses the model itself to generate feedback, guided by a
set of natural-language principles (the ``constitution'').  It has two
phases.

\textbf{SL-CAI:} The model generates an initial response, then critiques
that response against specific constitutional principles (``Is this
response harmful?''), then revises the response to address the critique.
The revised responses become SFT training data.

\textbf{RL-CAI (RLAIF):} The model generates preference labels (which of
two responses better satisfies the constitution?), and these AI-generated
preferences are used to train a reward model for standard RLHF.

CAI works because: (1)~it leverages the model's existing instruction-following
ability for self-evaluation, an ability that improves with scale;
(2)~the constitution provides compositional, auditable control over behavior
(unlike a black-box RM); (3)~AI feedback can be generated at any scale and
iterated rapidly, removing the human annotation bottleneck.  The surprising
finding is that RLAIF produces alignment quality comparable to human-feedback
RLHF, suggesting that the model can serve as a reasonable proxy for human
judgment when guided by explicit principles.

\redflags
\begin{itemize}
\item Confuses CAI with general RLHF.
\item Does not mention the constitution (the set of principles).
\item Cannot explain the SL-CAI vs.\ RL-CAI distinction.
\item Dismisses self-critique as ``the model grading its own homework.''
\end{itemize}

\interviewtest{Understanding of scalable oversight and whether the candidate
recognizes the tension between the expense of human feedback and the need
for alignment at scale.}

\goodgreat
\begin{itemize}
\item \badgeLfive{L5}: Describes both phases and the role of the
  constitution.  Explains why it reduces dependence on human annotations.
\item \badgeLsix{L6}: Discusses the quality of RLAIF vs.\ RLHF,
  the compositionality advantage of principles over a monolithic RM,
  and how the constitution can be updated without retraining the RM.
\item \badgeLseven{L7}: Connects to the broader scalable oversight
  problem, discusses limitations (model may fail to catch its own
  blind spots), and considers recursive reward modeling and debate as
  alternative approaches.
\end{itemize}

\followups{How does RLAIF compare to RLHF in quality?  What are the
limitations of self-critique?  What is the scalable oversight problem?}
\end{interviewq}

% --- Question 5 ---
\begin{interviewq}
\textbf{Q: Compare RLHF, DPO, and KTO.  When would you use each?}

\badgeLsix{L6} \quad \badgetype{Trade-off} \quad
\badgefreq{\freqmed}

\stronganswer
\textbf{RLHF (PPO):} The most general approach.  Requires a separate
reward model and RL optimization.  Highest complexity (4 models) but
provides an explicit reward signal for monitoring, enables online
exploration, and gives the most control over the optimization process.
Best for frontier labs with significant engineering resources where maximum
alignment quality matters and the ability to monitor reward during training
is valuable.

\textbf{DPO:} Collapses the reward model and RL into a single supervised
loss on pairwise preferences.  Simpler (2 models), more stable, lower
compute.  Best default choice for most teams doing alignment---gets 90--95\%
of RLHF quality with a fraction of the complexity.  Limited by its offline
nature and lack of explicit reward monitoring.

\textbf{KTO:} Only requires pointwise binary feedback (good/bad), not
pairwise comparisons.  This is a major practical advantage: you can
collect this data from production signals (user thumbs-up/down, regeneration
clicks) rather than constructing explicit comparison tasks.  Best when
pairwise preference data is unavailable or expensive to collect, but you
have abundant implicit feedback signals.

Decision framework: start with DPO (simplest effective method); move to
RLHF if you need tighter control or online exploration; use KTO when you
only have binary feedback data.

\redflags
\begin{itemize}
\item Cannot articulate the data requirements of each method.
\item Says ``DPO is always better because it is simpler.''
\item Does not know what distinguishes KTO from DPO (pointwise vs.\ pairwise).
\item Cannot describe a scenario where RLHF's complexity is justified.
\end{itemize}

\interviewtest{Whether the candidate can make nuanced engineering tradeoff
decisions rather than defaulting to a single approach.  This tests practical
judgment about alignment methods.}

\goodgreat
\begin{itemize}
\item \badgeLfive{L5}: Correctly describes all three methods, their data
  requirements, and when to use each.
\item \badgeLsix{L6}: Discusses the online vs.\ offline distinction in
  depth, mentions iterative DPO as a middle ground, and considers the
  monitoring advantage of having an explicit reward model.
\item \badgeLseven{L7}: Discusses the theoretical differences in
  optimization landscape (DPO can have degenerate solutions), the
  relationship between data scale and method choice, and how hybrid
  approaches (DPO + rejection sampling) combine benefits.
\end{itemize}

\followups{What is iterative DPO?  When does RLHF's online nature matter
most?  How would you collect KTO-style data from a production system?}
\end{interviewq}

% --- Question 6 ---
\begin{interviewq}
\textbf{Q: What makes a good reward model?  How do you evaluate it?}

\badgeLfive{L5} \quad \badgetype{System Design} \quad
\badgefreq{\freqmed}

\stronganswer
A good reward model should: (1)~have high \textbf{agreement} with human
preferences on a held-out test set (accuracy of 65--75\% is typical, given
that human inter-annotator agreement is only 70--80\%); (2)~\textbf{generalize}
beyond its training distribution to novel prompts and response styles;
(3)~be \textbf{robust} to surface-level features (length, formatting) that
correlate with but do not cause quality; (4)~provide \textbf{calibrated}
scores---the margin between good and bad responses should be meaningful.

Evaluation metrics: (a)~pairwise accuracy on held-out preferences;
(b)~correlation between RM score and human evaluation on a diverse test set;
(c)~sensitivity analysis---does the RM prefer longer responses regardless
of content? Does it prefer sycophantic responses?  (d)~distribution of
scores---are scores well-spread or collapsed to a narrow range?
(e)~performance on adversarial examples designed to test specific failure
modes (verbosity, sycophancy, formatting bias).

Practical improvements: train on diverse preference data across many topics
and difficulty levels; use RM ensembles for robustness; periodically retrain
on preferences from the \emph{current} policy (not just the SFT model);
consider larger RM models (RM quality scales with model size).

\redflags
\begin{itemize}
\item Only mentions accuracy without discussing generalization or robustness.
\item Does not consider systematic biases (length, formatting).
\item Cannot describe how to evaluate an RM beyond ``test accuracy.''
\item Ignores the relationship between RM quality and downstream RLHF quality.
\end{itemize}

\interviewtest{Whether the candidate understands that the RM is the
bottleneck of alignment and can think critically about evaluation beyond
simple metrics.}

\goodgreat
\begin{itemize}
\item \badgeLfive{L5}: Lists the key properties (accuracy, generalization,
  robustness) with concrete evaluation strategies.
\item \badgeLsix{L6}: Discusses bias analysis (length, sycophancy,
  formatting), RM ensembles, the upper bound set by inter-annotator
  agreement, and how RM errors propagate through PPO.
\item \badgeLseven{L7}: Discusses reward model overoptimization bounds,
  the relationship between RM size and hackability, and how process
  reward models provide a different evaluation paradigm.
\end{itemize}

\followups{How does inter-annotator agreement limit RM quality?  How
would you detect length bias in a reward model?  What is the effect of
RM size on alignment quality?}
\end{interviewq}

% --- Question 7 ---
\begin{interviewq}
\textbf{Q: Derive the DPO loss from the RLHF objective.}

\badgeLseven{L7} \quad \badgetype{Mathematical} \quad
\badgefreq{\freqmed}

\stronganswer
Start with the KL-constrained objective:
$\max_\pi \E_{y \sim \pi}[r(x,y)] - \beta \KL(\pi \| \pi_{\text{ref}})$.

\textbf{Step 1:} Find the optimal policy.  Expanding the KL and taking the
functional derivative with respect to $\pi(y|x)$:
$r(x,y) - \beta \log(\pi(y|x)/\pi_{\text{ref}}(y|x)) - \beta = \lambda$
(Lagrange multiplier for normalization).  Solving:
$\pi^*(y|x) = \frac{1}{Z(x)} \pi_{\text{ref}}(y|x) \exp(r(x,y)/\beta)$
where $Z(x)$ is the partition function.

\textbf{Step 2:} Reparameterize the reward.  Take log:
$r(x,y) = \beta \log(\pi^*(y|x)/\pi_{\text{ref}}(y|x)) + \beta \log Z(x)$.

\textbf{Step 3:} Substitute into the Bradley-Terry preference model.  The
RM loss uses $\sigma(r(x,y_w) - r(x,y_l))$.  Computing the reward
difference: the $\beta \log Z(x)$ terms cancel (both are conditioned on
the same $x$), giving:
$r(x,y_w) - r(x,y_l) = \beta[\log(\pi^*(y_w|x)/\pi_{\text{ref}}(y_w|x)) -
\log(\pi^*(y_l|x)/\pi_{\text{ref}}(y_l|x))]$.

\textbf{Step 4:} Replace $\pi^*$ with the trainable policy $\pi_\theta$:
$\loss_{\text{DPO}} = -\E[\log \sigma(\beta(\log \frac{\pi_\theta(y_w|x)}{\pi_{\text{ref}}(y_w|x)} - \log \frac{\pi_\theta(y_l|x)}{\pi_{\text{ref}}(y_l|x)}))]$.

The critical steps are: the closed-form optimal policy (Step~1), the
reward reparameterization (Step~2), and the cancellation of the partition
function in the reward difference (Step~3).

\redflags
\begin{itemize}
\item Cannot state the KL-constrained objective as the starting point.
\item Does not derive the optimal policy (just states the DPO loss).
\item Misses the partition function cancellation---the key trick.
\item Cannot explain why replacing $\pi^*$ with $\pi_\theta$ is valid.
\end{itemize}

\interviewtest{Deep mathematical understanding of DPO.  This is a
principal-level question that tests whether the candidate can derive
results from first principles rather than memorizing final formulas.}

\goodgreat
\begin{itemize}
\item \badgeLfive{L5}: States the DPO loss and explains at a high level
  that it comes from reparameterizing the reward.
\item \badgeLsix{L6}: Walks through the four steps above with the key
  equations, correctly identifies where $Z(x)$ cancels.
\item \badgeLseven{L7}: Completes the full derivation including the
  Lagrangian, discusses the relationship between $\beta$ in DPO and in
  RLHF, analyzes the gradient of the DPO loss (which tokens get the
  strongest gradient signal), and considers conditions under which the
  DPO objective has non-unique solutions.
\end{itemize}

\followups{What does the gradient of the DPO loss look like?  Which
tokens receive the strongest learning signal?  Under what conditions
can DPO fail to converge to the optimal RLHF policy?}
\end{interviewq}

% --- Question 8 ---
\begin{interviewq}
\textbf{Q: What is the difference between process and outcome reward
models?  When would you use each?}

\badgeLsix{L6} \quad \badgetype{Conceptual} \quad
\badgefreq{\freqmed}

\stronganswer
An \textbf{outcome reward model} (ORM) evaluates the complete response as a
whole, assigning a single scalar reward to the entire output.  A
\textbf{process reward model} (PRM) evaluates each step of a reasoning
chain individually, scoring whether each step is correct given the
preceding steps.

The key advantage of PRMs is \textbf{credit assignment}: when a math
solution is wrong, the ORM says ``this is bad'' without indicating which
step failed.  The PRM pinpoints the exact error, providing denser and more
informative supervision.  Lightman et al.\ (2023) showed that PRMs
significantly outperform ORMs for mathematical reasoning tasks.

PRMs also enable \textbf{search over reasoning paths}: the model generates
multiple candidate next steps, the PRM scores each, and search algorithms
(beam search, MCTS) select the most promising paths.  This is the
mechanism behind test-time compute scaling in models like o1.

Use an ORM for general-purpose alignment (chat, instruction following)
where there is no natural decomposition into steps.  Use a PRM when the
task has verifiable intermediate steps (mathematics, multi-step
reasoning, code generation) and when you want to enable test-time search.
The main limitation of PRMs is data cost: step-level annotations are
5--10$\times$ more expensive than response-level annotations, though
automated verification (e.g., executing code to check intermediate results)
can reduce this.

\redflags
\begin{itemize}
\item Cannot define the difference between process and outcome supervision.
\item Does not mention credit assignment as the key advantage.
\item Thinks PRMs are always better without acknowledging the data cost.
\item Cannot explain how PRMs enable test-time search.
\end{itemize}

\interviewtest{Awareness of cutting-edge alignment research (PRMs, test-time
compute) and the ability to reason about when fine-grained supervision is
worth the added cost.}

\goodgreat
\begin{itemize}
\item \badgeLfive{L5}: Correctly distinguishes ORM and PRM, mentions
  credit assignment and the data cost tradeoff.
\item \badgeLsix{L6}: Discusses Lightman et al.'s results, explains
  how PRMs enable tree search (beam search, MCTS over reasoning steps),
  and considers automated process supervision.
\item \badgeLseven{L7}: Connects to test-time compute scaling (o1-style
  reasoning), discusses the relationship between process supervision and
  formal verification, and considers how PRMs change the scaling laws
  (trading training compute for inference compute).
\end{itemize}

\followups{How does test-time compute scaling work?  How would you
collect process supervision data without human annotators?  What is the
relationship between PRMs and MCTS?}
\end{interviewq}

\bigskip
\begin{keyinsight}[Chapter Summary---What Interviewers Want]
This chapter covered the techniques that transform powerful but unaligned
base models into useful assistants.  The three most critical takeaways are:
(1)~the \textbf{RLHF pipeline} (SFT $\to$ RM $\to$ PPO)---know every
stage, why it exists, and what can go wrong; (2)~\textbf{DPO}'s
mathematical derivation---this tests whether you understand alignment
at a principled level, not just as a recipe; and (3)~\textbf{reward
hacking} as the central failure mode---Goodhart's law applied to
alignment is the deepest conceptual point in this chapter, and candidates
who can discuss it with sophistication signal genuine understanding of
why alignment is hard.  Interviewers at frontier labs increasingly test
these topics, and candidates who treat alignment as ``just fine-tuning''
reveal a gap in their understanding of modern NLP.
\end{keyinsight}
